\section{Admissible classes for integer certificates}
\label{sec:lattice}

For spatial certificates, the permitted pieces constrain how the domain can
be covered. Integer certificates have a related distinction: a collection of
convex pieces may cover the required integer points without covering all
fractional points of the root relaxation. This section identifies the least
number of pieces in that semantic model, then separates it from the cost of
using variable or split disjunctions. The same distinction explains why an
integer-space cut can shorten a particular tree without improving the lower
bound forced by the projected relaxation.

\subsection{The projected relaxation and the required integer points}

Consider
\begin{equation}
 \OPT=\inf\{f(x,y):(x,y)\in C,\ x\in\Z^n\},
 \label{lattice:problem}
\end{equation}
where $y$ denotes the continuous variables. Let $F\subseteq\Z^n$ be the
feasible integer projections, and let $f_F(x)$ be the infimum of the objective
over the feasible continuous variables with integer part $x$. A fixed convex
relaxation consists of a convex set $\widehat K\supseteq C$ and a jointly
convex objective $\widehat\phi$ that underestimates $f$ on $C$. Its projected
value function is
\begin{equation}
 \phi(x)=\inf\{\widehat\phi(x,y):(x,y)\in\widehat K\},
 \qquad r(Q)=\inf_{x\in Q}\phi(x).
 \label{lattice:projection}
\end{equation}
We assume that $\phi:\R^n\to\R\cup\{+\infty\}$ is proper. Infimal
projection preserves convexity, although it need not preserve lower
semicontinuity. The convention is $\inf\varnothing=+\infty$; in particular,
$\phi=+\infty$ outside its effective domain $K$. The value $r(Q)$ is the node
bound when branching restricts only the integer variables to $Q$.

Fix a set $P$ with $F\subseteq P\subseteq\Z^n$. It records which integer
points the certificate must preserve. Variable and split disjunctions on a
bounded root box usually take $P$ to be all integer points in that box;
branching justified only on feasible points may take $P=F$. These choices
can give different lower bounds. Fix also a threshold $\tau\in\R\cup
\{+\infty\}$. An additive $\varepsilon$ certificate uses
$\tau=\OPT-\varepsilon$, for finite $\OPT$ and $\varepsilon\ge0$;
infeasibility uses $\tau=+\infty$.

\begin{definition}[$P$-covering convex-piece certificate]
\label{lattice:tree-definition}
A finite rooted tree assigns a convex set $Q_v\subseteq\R^n$ to each node.
The root contains $P$. Each internal node $v$ satisfies
\begin{equation}
 Q_v\cap P\ \subseteq\ \bigcup_{w\text{ child of }v}Q_w.
 \label{lattice:integer-cover}
\end{equation}
Its effective set is $\overline Q_v=\bigcap_{u\preceq v}Q_u$, where the
intersection is over its root path. The tree is a $\tau$ certificate if
$r(\overline Q_\ell)\ge\tau$ at every leaf $\ell$.
\end{definition}

The sets need not be closed. Replacing every node set by its effective set
makes the tree nested and preserves \eqref{lattice:integer-cover}: a required
point in $\overline Q_v$ belongs to at least one child and to all its
ancestors. The children must cover only $P\cap Q_v$. No coverage of all
fractional points is required. This is the semantic certificate class used
below. Variable branches $x_i\le c$ or $x_i\ge c+1$, integer splits
$\pi^Tx\le c$ or $\pi^Tx\ge c+1$, and finite multiway integer hyperplane
branches are subclasses when their integer disjunctions cover $P$.

A completed branch-and-bound run gives such a certificate if it preserves
$P$ and prunes only by relaxation infeasibility, a valid computed bound
$\beta_v\le r(\overline Q_v)$ satisfying $\beta_v\ge UB_v-\varepsilon$, or
an attained relaxation optimum that is feasible for
\eqref{lattice:problem} and has the same original and relaxed objective.
Indeed, every feasible incumbent satisfies $UB_v\ge\OPT$. Bound pruning
therefore gives $r(\overline Q_v)\ge\OPT-\varepsilon$, and exact integral
pruning gives $r(\overline Q_v)\ge\OPT$. This argument permits weaker
computed lower bounds and arbitrary node selection. Reductions that remove
required integer points need the separate accounting below.

\subsection{Class number and exact semantic attainment}

\begin{definition}[Admissibility and class number]
\label{lattice:class-definition}
A set $I\subseteq P$ is $\tau$-admissible if
\begin{equation}
 r(\operatorname{conv}I)\ge\tau.
 \label{lattice:admissibility}
\end{equation}
The class number $\kappa_\tau(P)$ is the least number of admissible sets
covering $P$, or $+\infty$ if no finite cover exists. Set
$\kappa_\tau(\varnothing)=0$.
\end{definition}

Admissibility is hereditary. A cover may therefore be turned into a
partition by assigning every point to its first covering class. Pairwise
conflicts provide useful lower bounds, but they can miss obstructions that
involve more than two points. Define the midpoint graph on $P$ by
$a\sim b$ when $a\ne b$ and $\phi((a+b)/2)<\tau$. Define the segment graph
by $a\sim b$ when $\inf_{[a,b]}\phi<\tau$. An admissible class is an
independent set in both graphs. Thus
\begin{equation}
 \omega(G^{\rm mid}_\tau)\le\omega(G^{\rm seg}_\tau)
 \le\chi(G^{\rm seg}_\tau)\le\kappa_\tau(P).
 \label{lattice:graph-bound}
\end{equation}
For every nonempty finite $S\subseteq P$ with positive denominator, there
is also the counting bound
\begin{equation}
 \kappa_\tau(P)\ge
 \frac{|S|}{\max\{|I\cap S|:I\text{ is admissible}\}}.
 \label{lattice:count-bound}
\end{equation}
These follow by restricting an admissible partition to pairs and to $S$,
respectively. If no admissible class contains a point of $S$, no admissible
cover exists and $\kappa_\tau(P)=+\infty$. Midpoint conflicts and bounded
occupancy of a leaf extend established branch-and-bound lower-bound
mechanisms \citep{deyDubeyMolinaro2023BranchAndBoundLowerBounds}. Hiding
sets give a related obstruction to small formulations
\citep{kaibelWeltge2015LowerBoundsBranchAndBound}.

\begin{theorem}[Class number and semantic trees]
\label{lattice:class-number}
Every $P$-covering convex-piece $\tau$ certificate has at least
$\kappa_\tau(P)$ leaves. If $P\ne\varnothing$ and
$k=\kappa_\tau(P)<\infty$, the semantic model of
\cref{lattice:tree-definition} attains this lower bound with $k$ leaves.
It also attains it with a full binary tree having $2k-1$ nodes. If $P$ is
finite, all node sets in an attaining tree can be polytopes.
\end{theorem}

\begin{proof}[Proof of the lower bound]
Route each $p\in P$ from the root to a child containing it and continue
until a leaf is reached. Let $I_\ell$ consist of the points routed to leaf
$\ell$. Then $I_\ell\subseteq\overline Q_\ell$, and convexity gives
$\operatorname{conv}I_\ell\subseteq\overline Q_\ell$. Consequently
$r(\operatorname{conv}I_\ell)\ge r(\overline Q_\ell)\ge\tau$. The nonempty
$I_\ell$ form an admissible partition of $P$.
\end{proof}

The attainment proof in \cref{app:lattice:semantic} uses separation from the
convex strict sublevel set $E_\tau=\{\phi<\tau\}$. Each class can be
enlarged to a \emph{hemispace}: a convex set whose complement is convex and
which misses $E_\tau$. Successively removing these hemispaces gives a
binary chain. At its last internal node, the two remaining children cover
the remaining points of $P$; they need not cover its fractional points.
The appendix proves the required hemispace separation for arbitrary
disjoint convex sets, including nonclosed sets.

The root is part of the semantic certificate: it may be any convex set
containing $P$. This convention matters when $k=1$, since
$r(\operatorname{conv}P)\ge\tau$ can hold while a prescribed larger
continuous root has bound below $\tau$. For $k\ge2$ the binary construction
can begin at any prescribed convex root containing $P$. For $k=1$, its
one-node form requires the permitted integer-preserving root restriction.
The lower bound holds without this extra permission.

Exact attainment is an existence statement about this specified certificate
class. A solver may require disjunctions whose coverage of all node integer
points has a short, explicit proof, and may restrict itself to variable or
split branching. Finding a minimum admissible partition, separating its
classes, and checking its integer coverage can be costly. Finite $P$ makes
the node sets polyhedral, but does not bound the number of vertices or the
time required to construct and verify them. This separates certificate size
from construction and checking costs, as in the continuous distinction
between covers, rectangular partitions, and permitted trees.

\subsection{Cuts and certified integer removals}

\begin{theorem}[Preserved classes and removal accounting]
\label{lattice:operations}
Suppose all node bounds use the fixed projected value function $\phi$.
\begin{enumerate}[label=(\alph*)]
\item Convex cuts in the integer variables that contain all required integer
points at the node preserve the leaf lower bound
$L\ge\kappa_\tau(P)$. In particular, cuts valid for all feasible integer
points preserve $L\ge\kappa_\tau(F)$, even if the cuts give the node's
integer hull. Convex feasibility reductions preserving those points have
the same property.
\item Suppose a node's effective set undergoes a finite nested sequence
$Q^{(0)}\supseteq\cdots\supseteq Q^{(s)}$, and every removal has a convex
test set $D_i$ satisfying
\begin{equation}
 Q^{(i-1)}\cap P\subseteq Q^{(i)}\cup D_i,
 \qquad r(Q^{(i-1)}\cap D_i)\ge\tau.
 \label{lattice:removal-hypothesis}
\end{equation}
If the remaining tree is pruned at threshold $\tau$, then
\begin{equation}
 \kappa_\tau(P)\le L+S,
 \qquad S=\sum_v s_v.
 \label{lattice:augmented-count}
\end{equation}
The same bound holds for simultaneous certified pieces covering the removed
integer points, each tested on the unreduced effective set.
\item For box-bound changes, one removal is one change of one integer bound
by any amount. One pass changing each bound at most once per node gives
$N\ge\kappa_\tau(P)/(2n+1)$. One such pass only at the root gives
$L\ge\kappa_\tau(P)-2n$. For binary variables, iterated rounded bound
changes give $N\ge\kappa_\tau(P)/(n+1)$.
\end{enumerate}
\end{theorem}

The proof in \cref{app:lattice:operations} assigns the removed integer
points to additional certified leaves. For general integer variables,
iterated passes require the actual number $S$ of bound changes; the proof
does not give a factor depending only on dimension. Changes of the same
bound at different times need not combine into a single admissible class.

Integer rounding gives a clean removal certificate. For example, suppose
the current set is $Q$ and incumbent-based tightening finds
\begin{equation*}
 \ell'_i=\left\lceil\inf\{x_i:x\in Q,\ \phi(x)\le c\}\right\rceil,
 \qquad c=UB-\varepsilon\ge\tau,
\end{equation*}
with a finite infimum. Every point of
$Q\cap\{x_i\le\ell'_i-1\}$ has $\phi(x)>c$, so its infimum is at least
$c$. Together with the retained set $Q\cap\{x_i\ge\ell'_i\}$, this closed
slab covers the integer points of $Q$. Its endpoint is one integer unit
below the retained bound. Thus the retained-boundary obstruction for
continuous closed discarded slabs does not arise. Upper-bound rounding is
analogous. Reduced-cost certificates and incumbent-based probing fit
\eqref{lattice:removal-hypothesis} when their tested pieces have the stated
bound on the current node set.

The leaf bound can fail without the removal count. Take the root
$[0,1]$, $P=\{0,1\}$, and $\phi(x)=(x-0.4)^2$ on that box. The integer
optimum is $0.16$. At $\varepsilon=0.01$, integers $0$ and $1$ conflict,
so the class number is two. Tightening with the closed incumbent cutoff
$\phi\le UB=0.16$ gives $x\le0.8$; integer rounding retains only
$\{0\}$. This permits a one-node certificate, and the removed point
$\{1\}$ is the second certificate piece. Strengthening the relaxation in continuous or epigraph
variables changes $\phi$ and can change its class number. Symmetry or
dominance reductions without a certificate of the form
\eqref{lattice:removal-hypothesis} are outside this accounting.

\subsection{The cost of permitted disjunctions}

Two elementary comparisons show why $\kappa_\tau$ cannot by itself predict
the minimum tree size for a restricted branch class.

\begin{proposition}[Linear descriptions and projected descriptions]
\label{lattice:linear}
For a pure integer linear program with natural relaxation
$K=\{x:Ax\le b\}$, integral $A,b$ with $m$ rows, and objective $c^Tx$,
$\kappa_\tau(\Z^n)\le m+1$ for every finite $\tau\le\OPT$. If it is
infeasible, $\kappa_{+\infty}(\Z^n)\le m$. For a rational mixed-integer
linear relaxation, if $\tau<\OPT$ and the projected closed sublevel
polyhedron $\{x:\exists y,\ (x,y)\in\widehat K,
\ \widehat\phi(x,y)\le\tau\}$ has a rational description with $q$ rows,
then $\kappa_\tau(\Z^n)\le\max(q,1)$.
\end{proposition}

\begin{proof}
In the pure case, use the integer points in
$\{c^Tx\ge\OPT\}$ and in $\{a_j^Tx\ge b_j+1\}$, $1\le j\le m$.
They cover $\Z^n$: an integer point with objective below $\OPT$ violates
some row by at least one. Their hulls respectively have objective at least
$\OPT$ or miss $K$. In the infeasible case, only the violated-row classes
are needed. In the mixed case, scale the $q$ projected rows to integral
left-hand sides $g_j^Tx\le h_j$. The integer-free sublevel polyhedron is
avoided by the covering halfspaces
$g_j^Tx\ge\lfloor h_j\rfloor+1$. On each halfspace the projected value is
at least $\tau$, since every feasible continuous extension has value
strictly above $\tau$. If the sublevel polyhedron is empty, one class
suffices.
\end{proof}

Projection can hide exponentially many classes in a short formulation.
The mixed-integer linear problem
\begin{equation*}
 \min\sum_{i=1}^n s_i,
 \qquad s_i\ge x_i-\tfrac12,\quad s_i\ge\tfrac12-x_i,
 \qquad x\in\Z^n
\end{equation*}
has $2n$ rows and projected function
$\phi(x)=\|x-\tfrac12\mathbf1\|_1$. Its optimum is $n/2$. For
$0\le\varepsilon<1/2$, all $2^n$ binary vectors pairwise conflict: vectors
at Hamming distance $d$ have midpoint value $(n-d)/2$.
Conversely, the $2^n$ halfspaces
$\sigma^T(x-\tfrac12\mathbf1)\ge n/2$, $\sigma\in\{-1,1\}^n$, cover
$\Z^n$ and are admissible. Hence $\kappa_{\OPT-\varepsilon}(\Z^n)=2^n$.

\begin{proposition}[A quadratic separation for variable branching]
\label{lattice:jeroslow}
Let $n$ be odd, $h=(n+1)/2$, and minimize
$\phi(x)=(\mathbf1^Tx-n/2)^2$ over $\{0,1\}^n$, with the natural box
relaxation. For $0\le\varepsilon<1/4$, the optimum is $1/4$ and the class
number is two. The split
$\mathbf1^Tx\le h-1$ or $\mathbf1^Tx\ge h$ certifies the optimum with
three nodes. Every variable-branching certificate has at least
$\binom{n+1}{h}$ leaves. This count is attained when each branch fixes a
free binary variable and every node whose bound reaches the threshold is
immediately made a leaf.
\end{proposition}

\begin{proof}
The two sum halfspaces cover the integer points and have bound $1/4$;
the convex hull of all binary points has bound zero, so one class cannot
suffice. A variable node fixing $a$ ones and $b$ zeros has bound zero
precisely when $a,b\le h-1$. Otherwise its bound is at least $1/4$.
Contract redundant branches and make every already-certified node a leaf.
At every remaining internal node a free variable is fixed, giving the two
states $(a+1,b)$ and $(a,b+1)$. The resulting leaf count satisfies
$T(a,b)=T(a+1,b)+T(a,b+1)$ for $a,b<h$, with boundary value one at
$\max(a,b)=h$. Pascal's identity gives
$T(a,b)=\binom{2h-a-b}{h-a}$, hence $T(0,0)=\binom{n+1}{h}$.
\end{proof}

This is the parity obstruction in a convex quadratic optimization form.
It separates variable branching from both semantic branching and integer
splits. It does not establish a comparable separation between general
split trees and semantic trees for one convex quadratic.

\begin{theorem}[Lattice-free interpretation]
\label{lattice:facet-number}
Let $\phi:\R^n\to\R$ be convex and
$\tau\le\inf_{z\in\Z^n}\phi(z)$. If $E_\tau=\{\phi<\tau\}$ is empty,
then $\kappa_\tau(\Z^n)=1$. Otherwise it equals the least number of closed
halfspaces covering $\Z^n$ and avoiding $E_\tau$, and the least number of
facets of a polyhedron $M$ satisfying
\begin{equation}
 E_\tau\subseteq\operatorname{int}M,
 \qquad\operatorname{int}M\cap\Z^n=\varnothing.
 \label{lattice:enclosing-polyhedron}
\end{equation}
In particular, $\kappa_\tau(\Z^n)\le2^n$.
\end{theorem}

The proof is in \cref{app:lattice:facets}. The final bound uses the standard
maximal lattice-free polyhedron theorem
\citep{basuConfortiCornuejolsZambelli2010MaximalLatticeFree}. This interpretation connects
class number to relative relaxation complexity: the object to be enclosed
is now a strict objective sublevel set. It retains the distinction between
the number of separating pieces and the cost of realizing them through
permitted integer disjunctions.

\subsection{Haar-random closest-vector certificates}

Let $L$ be a unimodular lattice drawn from Haar probability on
$\operatorname{SL}_n(\R)/\operatorname{SL}_n(\Z)$, $n\ge2$. Conditional
on $L$, let $t$ be uniform on the quotient torus $\R^n/L$. For any lattice
basis $B$, use the continuous relaxation
\begin{equation}
 \phi(x)=\|Bx-t\|^2,
 \qquad P=\Z^n,
 \qquad\OPT=\dist(t,L)^2.
 \label{lattice:cvp-model}
\end{equation}
Let $g_n$ be the radius of a Euclidean ball of volume one, so
$g_n^2\sim n/(2\pi e)$, and let $\lambda_1(L)$ be the shortest nonzero
lattice-vector length. A unimodular change of basis maps integer points and
convex certificate pieces bijectively. All class-number bounds below are
therefore independent of the chosen basis. This model is distinct from a
matrix with independent Gaussian columns scaled to determinant one.

\begin{lemma}[A deterministic occupancy bound]
\label{lattice:cap}
Suppose $\tau>0$ and $R^2\le\tau+\lambda_1(L)^2/2$. Every admissible
class contains at most $n+1$ lattice points $Bz$ in the open ball
$B^\circ(t,R)$. Consequently
\begin{equation}
 \kappa_\tau(\Z^n)\ge
 \frac{\#(L\cap B^\circ(t,R))}{n+1}.
 \label{lattice:cap-count}
\end{equation}
\end{lemma}

An admissible class lies on the far side of a supporting hyperplane to the
ball of radius $\sqrt\tau$. Its intersection with $B^\circ(t,R)$ lies in
a ball of radius at most $\lambda_1/\sqrt2$. The vectors from that ball's
center to its lattice points have strictly negative pairwise inner
products; there are at most $n+1$ such vectors. The full proof is in
\cref{app:lattice:cap-proof}. This uses occupancy rather than only pairwise
conflicts between all sampled points.

\begin{theorem}[Random closest-vector class bounds]
\label{lattice:cvp}
In \eqref{lattice:cvp-model}, put $\tau=\OPT-\varepsilon$ and
$q=\varepsilon/g_n^2$. Fix $\delta\in(0,0.1)$ and suppose
\begin{equation*}
 0\le q<(1-\delta)^2,
 \qquad R_-^2=\tfrac32(1-\delta)^2-q>1.
\end{equation*}
Then, with probability at least
$1-\tfrac32(1-\delta)^n-4R_-^{-n}$,
\begin{equation}
 \kappa_\tau(\Z^n)\ge\frac{R_-^n}{2(n+1)}.
 \label{lattice:cvp-lower}
\end{equation}
For every $\gamma>0$, an upper bound holds with probability at least
$1-(1+\delta)^{-n}-e^{-\gamma n}$:
\begin{equation}
 \kappa_\tau(\Z^n)
 \le \bigl(e^{\gamma n}+C n^2\log(n+1)\bigr)
       \bigl(1+(1+\delta)^2-q\bigr)^{n/2},
 \label{lattice:cvp-upper}
\end{equation}
where $C$ is an absolute constant when $0\le q\le1/2$ and
$0<\delta<0.1$. If $\tau\le0$, one class suffices.

For any deterministic tolerance sequence $q_n\to q_0\in[0,1/2)$, these
bounds imply, with probability tending to one,
\begin{equation}
 2^{(\frac12\log_2(3/2-q_0)-o(1))n}
 \le\kappa_{\OPT-\varepsilon_n}(\Z^n)
 \le2^{(\frac12\log_2(2-q_0)+o(1))n}.
 \label{lattice:cvp-exponents}
\end{equation}
\end{theorem}

The proof in \cref{app:lattice:cvp-proof} combines torus unfolding with
Siegel's mean-value theorem in the normalization stated by
\citet{skenderi2021RandomLatticesSiegel}. Under the joint lattice-target law, the number
of lattice points in a translated set of volume $V$ has mean $V$ and
variance $V$. No Poisson independence assumption is needed. The upper
bound covers nearby lattice points individually and distant points by
halfspaces whose directions form a spherical cap cover.

At vanishing relative tolerance $q_n=o(1)$,
\eqref{lattice:cvp-exponents} gives lower exponent
$\tfrac12\log_2(3/2)=0.29248\ldots$ and upper exponent $1/2$.
For fixed positive $q_0$, both bounds depend on the tolerance. In
particular, the upper exponent is strictly below $1/2$. Thus an exponent
$1/2$ cannot hold uniformly for every fixed positive relative tolerance.

\begin{proposition}[A midpoint-clique lower bound]
\label{lattice:cvp-clique}
In \eqref{lattice:cvp-model}, fix $\delta\in(0,0.1)$,
$\varepsilon\ge0$, $q=\varepsilon/g_n^2$, and
$\tau=\OPT-\varepsilon$. Let $r_0^2=(1-\delta)^2-q$, take $1<\rho<4/3$ with
$\rho r_0^2>1$, and put $V=(\rho r_0^2)^{n/2}$. With probability at least
\begin{equation*}
 1-(1-\delta)^n-4/V
   -2n\left(\frac{4(\rho-1)}\rho\right)^{n/2},
\end{equation*}
the midpoint graph has a clique of size at least $V/4$. In particular, when
$q_n=o(1)$ its clique number is at least
$2^{(\frac12\log_2(4/3)-o(1))n}$ with probability tending to one.
\end{proposition}

The proof in \cref{app:lattice:clique-proof} removes one endpoint of each
nonconflicting pair from a ball of lattice points. Its exponential rate,
$0.20752\ldots$, is smaller than the occupancy rate in
\eqref{lattice:cvp-lower}. These two lower bounds use different information
about admissibility; neither establishes the exact exponent of the class
number. Every $\Z^n$-covering convex-piece tree inherits the class lower
bound, including variable branching in any reduced basis and general
integer splits. Certified incumbent reductions inherit
\eqref{lattice:augmented-count}; a claim about nodes requires its stated
removal budget. Enumeration that omits out-of-radius values also requires
the two certified integer tails, adding at most two leaves per internal
node and giving a node lower bound of one third of the class bound.

\subsection{Curvature conflicts and stronger block relaxations}

The following finite construction places all conflicting points in the
feasible integer set, and retains a unique optimum. It shows directly how
changing a projected relaxation can eliminate an exponential class
obstruction.

Use $k$ mutually orthogonal two-dimensional feature blocks. Block $j$ has
unit features $u_j,v_j$, response $y_j$, and ridge $\lambda>0$. There are
$2k$ binary indicators $z$, with $\mathbf1^Tz\le k$. The perspective
relaxation projects to
\begin{equation}
 g(z)=\sum_{j=1}^k g_j(z_{u_j},z_{v_j}),
 \qquad
 g_j(a,b)=y_j^T\left(I_2+
       \frac{a u_ju_j^T+b v_jv_j^T}{\lambda}\right)^{-1}y_j,
 \label{lattice:block-function}
\end{equation}
on $0\le z\le1$, $\mathbf1^Tz\le k$, and is $+\infty$ elsewhere.
For a binary support it equals the ordinary ridge objective on that
support. Define $g_{ab,j}=g_j(a,b)$ for $a,b\in\{0,1\}$, and
\begin{equation*}
 g_j^*=\min(g_{10,j},g_{01,j}),\quad
 \Delta_j=|g_{10,j}-g_{01,j}|,\quad
 d_j=g_j^*-g_j(1/2,1/2).
\end{equation*}

\begin{theorem}[Curvature conflicts with a unique optimum]
\label{lattice:curvature-gadget}
Suppose
\begin{equation}
 \max_j(g_j^*-g_{11,j})
   <\min_j(g_{00,j}-g_j^*),
 \qquad \sum_j\Delta_j+\varepsilon<\min_j d_j.
 \label{lattice:gadget-hypotheses}
\end{equation}
Then $\OPT=\sum_jg_j^*$. If every $\Delta_j>0$, the optimum is unique.
All $2^k$ one-feature-per-block supports pairwise conflict at
$\tau=\OPT-\varepsilon$. Every $F$-covering convex-piece certificate
using \eqref{lattice:block-function} has at least $2^k$ leaves, also with
arbitrary convex cuts valid for the node's feasible indicators. With
certified binary-bound reductions the node lower bound is
$2^k/(2k+1)$.

If each block relaxation is instead the closed convex hull of its four
mixed-integer objective epigraphs, the root value equals $\OPT$. For every
$k$, rational data satisfying \eqref{lattice:gadget-hypotheses}, with all
$\Delta_j>0$ and a fixed positive range of tolerances, can be chosen with
encoding length polynomial in $k$.
\end{theorem}

The proof and the rational family are in \cref{app:lattice:gadget}.
The stronger root relaxation changes the projected value function in the
epigraph variables. It therefore does not contradict preservation of
classes under cuts involving only the indicator variables. The conclusion
concerns these relaxation and operation classes, rather than a universal
algorithmic lower bound.

\subsection{A sufficient condition for one-path trees}

The hard-class mechanism has a useful counterpart for easy certificates.
Restrict the integer domain to a root box $\mathcal B=[\ell,u]$, with
integral endpoints, and in this subsection set $P=\mathcal B\cap\Z^n$.
Let $(x^\circ,y^\circ)$ be a feasible point of objective $U$, and suppose
$\phi(x^\circ)\ge U-\varepsilon$. Assume its incumbent value $U$ is
available whenever a node is processed. Write
\begin{equation}
 \begin{aligned}
 r(\mathcal B\cap\{x_i\le x_i^\circ-1\})&\ge U-\varepsilon
       &&\text{if }\ell_i<x_i^\circ,\\
 r(\mathcal B\cap\{x_i\ge x_i^\circ+1\})&\ge U-\varepsilon
       &&\text{if }x_i^\circ<u_i.
 \end{aligned}
 \tag{C1}\label{lattice:C1}
\end{equation}

\begin{lemma}[One-path certificate]
\label{lattice:path}
Under \eqref{lattice:C1}, any proper variable-branching run that prunes
nodes reaching the incumbent threshold before branching processes at most
$2D+1$ nodes, where $D\le\sum_i(u_i-\ell_i)$. For binaries, $D\le n$.
If $(x^\circ,y^\circ)$ is optimal, $\phi(x^\circ)=\OPT$, and every
nonempty wrong one-sided box has bound strictly greater than $\OPT$, the
same conclusion holds without an initial incumbent under best-bound
selection using the exact node values $r$, provided processing the
singleton $x=x^\circ$ obtains a feasible extension of value $\OPT$.
\end{lemma}

\begin{proof}
Every node not containing $x^\circ$ lies in a wrong one-sided root box.
Its bound is at least $U-\varepsilon$, by monotonicity, so it is pruned.
At each branched node exactly one child contains $x^\circ$. Hence the
branched nodes form a chain. Proper branching strictly shortens one integer
range, giving its length bound and the total count $1+2D$. If the
containing path reaches a singleton, its bound is at least
$U-\varepsilon$, so it can be pruned. In the
best-bound variant every node containing $x^\circ$ has bound at most
$\phi(x^\circ)=\OPT$, whereas every off-path node has strictly larger
bound. Until the optimum is found, best-bound selection follows the
containing path; at its singleton endpoint the stated processing
hypothesis supplies $\OPT$. All off-path nodes can then be pruned.
\end{proof}

Condition \eqref{lattice:C1} is a sufficient certificate criterion, rather
than a necessary condition for small trees. Probing each wrong one-sided
box with this incumbent fixes every integer variable to $x^\circ$; the
result uses one run node and up to $2n$ probe certificates. These
one-sided classes and $\{x^\circ\}$ also give
$\kappa_{\OPT-\varepsilon}(P)\le2n+1$: all their bounds are at least
$U-\varepsilon\ge\OPT-\varepsilon$.

Even within this regime, the class number can be much smaller than the
variable tree. For $n\ge2$, take $\phi(z)=\sum_i(z_i-a)^2$ on
$[0,1]^n$, $a=1/(2n)$, and $\varepsilon=1/(8n^2)$.
The optimum is $1/(4n)$ at zero, and each wrong fixing has bound
$(1-a)^2\ge\OPT$. The two classes $\{0\}$ and
$\{z\in\{0,1\}^n:\mathbf1^Tz\ge1\}$ are admissible: the latter hull's
minimum is attained at $\mathbf1/n$ and equals $\OPT$. The root bound is
zero, so the class number is exactly two. A node fixing only $s$ zeros has
bound $s/(4n^2)$ and is prunable only when $s=n$. The path to zero therefore
has $n$ branches and $n+1$ leaves. This gives $2n+1$ nodes despite a
two-class semantic certificate.
